\chapter{Mathematical and Cryptological Synthesis}
\label{ch:mathematical-cryptological-synthesis}

This chapter gives the central mathematical reading of the machine-checked
HAETAE security proof.  Its purpose is to make the checked implication legible to a
cryptologist: the game experiment, the exact bound expression, the hardness
ledger, the vacuity checks, the local non-vacuity evidence, and the remaining
trusted base are stated in mathematical form and connected to their EasyCrypt
names.

The proof is conditional.  EasyCrypt checks the game transformations,
random-oracle programming, sampler freshness arguments, active-signing lifting,
and algebraic composition of loss terms.  It does not prove MLWE hardness,
Module-SIS hardness, or implementation fidelity outside the formal model.  The
correct theorem is therefore a theorem about the modeled HAETAE random-oracle
scheme.

\section{Formal security object}
\label{sec:synthesis-formal-object-final}

Let \(A\) be a chosen-message adversary with access to a random oracle \(H\) and
a signing oracle \(O\).  The random-oracle EUF-CMA experiment samples
\[
  (pk,sk) \leftarrow \mathsf{KeyGen}^{H}(),
\]
answers signing queries by \(\mathsf{Sign}^{H}(sk,\cdot)\), answers hash queries
by lazy sampling, and receives a final candidate
\[
  (m^\star,ctx^\star,\sigma^\star)\leftarrow A^{H,O}(pk).
\]
The success event is
\[
  \mathsf{Forge}
  =
  \left[
    \mathsf{Verify}^{H}(pk,m^\star,ctx^\star,\sigma^\star)=1
    \wedge
    (m^\star,ctx^\star)\notin Q_S
  \right],
\]
where \(Q_S\) is the signing-query log.  For the informal target HAETAE scheme,
the advantage is
\[
  \Adv_{\haetae,\rom}^{\eufcma}(A)=\Pr[\mathsf{Forge}].
\]

In EasyCrypt, the generic experiment is supplied by \ec{Sig\_ROM.ec}; the
HAETAE instantiation is supplied by \ec{HAETAE\_Scheme.ec}; and the hop sequence
is supplied by \ec{HAETAE\_HopGames.ec} and \ec{HAETAE\_Security.ec}.  The
checked theorem surface has the implication shape
\[
  \mathcal{H}_{\mathsf{hop}}(A,B_{\mathsf{MLWE}},B_{\mathsf{bimodal}})
  \wedge
  \mathcal{H}_{\mathsf{hard}}(B_{\mathsf{MLWE}},B_{\mathsf{bimodal}})
  \Longrightarrow
  \Adv_{\mathsf{ECModel},\rom}^{\eufcma}(A)
  \le
  \ec{haetae\_euf\_counted\_rom\_bound}.
\]
The hypotheses are the checked hop inequality that reaches the MLWE and bimodal
self-target MSIS games, together with the declared hardness bounds for the
resulting reduction games.  The right-hand side is not a placeholder.  It is an
explicit EasyCrypt operation whose definition is normalized in
\ec{HAETAE\_Reductions.ec}.  The constants \(q_H=q_S=2\) belong to that bound
ledger and to the budgeted wrapper lemmas, not to the syntax of the source
\eufcma{} experiment in \ec{Sig\_ROM.ec}.

\section{Exact checked bound expression}
\label{sec:synthesis-exact-bound-expression}

The relevant EasyCrypt definitions are the following.

\begin{lstlisting}[language=EasyCrypt,caption={Checked bound definitions in \ec{HAETAE\_Reductions.ec}.},label={lst:haetae-reductions-bound-definitions}]
op mlwe_hardness_term : real.
op module_sis_hardness_term : real.
op bimodal_to_msis_reduction_loss_term : real = 0%r.

op bimodal_selftarget_msis_hardness_term =
  module_sis_hardness_term + bimodal_to_msis_reduction_loss_term.

op signature_query_count : real = 2%r.
op hash_query_count : real = 2%r.
op signature_query_budget_count : int = 2.
op hash_query_budget_count : int = 2.
op challenge_support_cardinality_lower_bound : real =
  288230376151711744%r.

op rom_hash_query_budget : real = hash_query_count.
op rom_signature_query_budget : real = signature_query_count.
op rom_total_query_budget =
  rom_hash_query_budget + rom_signature_query_budget + 1%r.

op counted_rom_collision_term =
  (rom_total_query_budget * rom_total_query_budget) /
    challenge_support_cardinality_lower_bound.

op counted_rom_prequery_term =
  (rom_signature_query_budget * (rom_hash_query_budget + 1%r)) /
    challenge_support_cardinality_lower_bound.

op counted_rom_min_entropy_term =
  (rom_signature_query_budget * (rom_hash_query_budget + 1%r)) /
    challenge_support_cardinality_lower_bound.

op counted_rom_programming_loss_term =
  counted_rom_collision_term +
  counted_rom_prequery_term +
  counted_rom_min_entropy_term.

op haetae_euf_counted_rom_bound =
  mlwe_hardness_term +
  bimodal_selftarget_msis_hardness_term +
  counted_rom_programming_loss_term.
\end{lstlisting}

Mathematically, let
\[
  q_H=\ec{hash\_query\_count}=2,
  \qquad
  q_S=\ec{signature\_query\_count}=2,
\]
and let
\[
  N=\ec{challenge\_support\_cardinality\_lower\_bound}
   =288230376151711744=2^{58}.
\]
Then
\[
  L_{\mathsf{coll}}=\frac{(q_H+q_S+1)^2}{N},
  \qquad
  L_{\mathsf{pre}}=\frac{q_S(q_H+1)}{N},
  \qquad
  L_{\mathsf{me}}=\frac{q_S(q_H+1)}{N}.
\]
For the current constants this gives
\[
  L_{\mathsf{ROM}}
  =L_{\mathsf{coll}}+L_{\mathsf{pre}}+L_{\mathsf{me}}
  =\frac{25+6+6}{2^{58}}
  =\frac{37}{2^{58}}.
\]

This concrete value should not be mistaken for a fully general EUF-CMA
query-bound theorem.  It is the concrete ROM contribution used by the theorem
family checked in \ec{HAETAE\_Security.ec}; the constants and arithmetic
normalizations are defined in \ec{HAETAE\_Reductions.ec}, whose public constants set
\[
  \ec{hash\_query\_budget\_count}=2,
  \qquad
  \ec{signature\_query\_budget\_count}=2.
\]
The reusable algebraic shape is
\[
  \frac{(q_H+q_S+1)^2}{N}
  +
  \frac{q_S(q_H+1)}{N}
  +
  \frac{q_S(q_H+1)}{N},
\]
but the machine-checked implication documented here instantiates that expression at
the two-hash-query, two-signing-query wrapper constants.  A future generalized theorem
should promote \(q_H\), \(q_S\), and \(N\) from fixed operations to symbolic
parameters and recheck all monotonicity and subunit obligations.

Under the hop and hardness hypotheses of the final EasyCrypt theorem family, the
expanded checked counted-ROM conclusion is therefore
\[
\boxed{
  \Adv_{\mathsf{ECModel},\rom}^{\eufcma}(A)
  \le
  L_{\mathsf{MLWE}}
  + L_{\mathsf{ModSIS}}
  + L_{\mathsf{bimodal}\to\mathsf{MSIS}}
  + \frac{37}{2^{58}}
}
\]
where
\[
  L_{\mathsf{MLWE}}=\ec{mlwe\_hardness\_term},
\]
\[
  L_{\mathsf{ModSIS}}=\ec{module\_sis\_hardness\_term},
\]
and
\[
  L_{\mathsf{bimodal}\to\mathsf{MSIS}}
  =\ec{bimodal\_to\_msis\_reduction\_loss\_term}.
\]
The current EasyCrypt file defines the last term as \(0\), but it is still
listed separately because it records a proof-design decision: bimodal
self-target MSIS is grouped with Module-SIS through
\ec{bimodal\_selftarget\_msis\_hardness\_term}.

For precision, this result should be cited as a checked fixed-ledger implication
for the formal EasyCrypt model, with the bound constants instantiated at
\(q_H=q_S=2\), conditional on the EasyCrypt antecedents of the selected
composition lemma.  It is not yet the conventional
all-polynomial-time, arbitrary-query EUF-CMA theorem.  The dissertation's
strongest checked claim is
\[
  \mathcal{H}_{\mathsf{EC}}
  \Longrightarrow
  \Adv_{\mathsf{ECModel},\rom}^{\eufcma}(A)
  \le
  \ec{haetae\_euf\_counted\_rom\_bound}.
\]
Here \(\mathcal{H}_{\mathsf{EC}}\) abbreviates the explicit antecedents of the
EasyCrypt composition lemma being cited: the final hop/reduction inequality and
the MLWE and Module-SIS hardness bounds for the declared reduction modules
supplied to the theorem.  The final composition lemma does not itself construct
those MLWE or Module-SIS solvers from \(A\); it composes the explicitly stated
probability and hardness premises.
One of those antecedents is currently a proof-surface premise rather than a
useful concrete-security premise: \ec{fs\_with\_aborts\_min\_entropy\_term} is
defined as \(12\), and the current sources prove
\(\ec{rejection\_sampling\_loss\_term}\ge 12\).  Therefore a final premise that
adds \(\ec{rejection\_sampling\_loss\_term}\) to a probability cannot be
satisfied by the sum of two game-success probabilities without further repair to
the premise or loss normalization.
Any statement for larger or symbolic budgets is a proposed generalization unless
the corresponding EasyCrypt constants and lifting lemmas are replaced by
parameterized versions and rechecked.

\begin{theorem}[Checked scoped counted-ROM implication]
\label{thm:checked-counted-rom-ledger}
For every adversary satisfying the declared EasyCrypt module interfaces, with
the checked proof and bound ledger instantiated through the fixed wrapper
constants \(q_H=q_S=2\), assume the final checked hop/reduction premise
\[
\begin{aligned}
  &\Pr[\mathsf{UF\_NMA}_{\mathsf{ROExactHyperballPaperSim}}(A)\Rightarrow 1]
  +\ec{rejection\_sampling\_loss\_term} \\
  &\qquad\le
  \Pr[\ec{MLWE\_Game}(B_{\mathsf{MLWE}})\Rightarrow 1]
  +
  \Pr[\ec{BimodalSelfTargetMSIS\_Game}(B_{\mathsf{bimodal}})\Rightarrow 1],
\end{aligned}
\]
and the hardness premises
\[
  \Pr[\ec{MLWE\_Game}(B_{\mathsf{MLWE}})\Rightarrow 1]
  \le \ec{mlwe\_hardness\_term}
\]
and
\[
  \Pr[\ec{ModuleSIS\_Game}
      (\ec{BimodalMSIS\_As\_ModuleSIS}(B_{\mathsf{bimodal}}))\Rightarrow 1]
  \le \ec{module\_sis\_hardness\_term}.
\]
Then the machine-checked implication establishes the probability-only
random-oracle-model conclusion for the formal EasyCrypt scheme model
\(\mathsf{ECModel}\):
\[
\begin{aligned}
  \Adv_{\mathsf{ECModel},\rom}^{\eufcma}(A)
  \le{}&
  \ec{mlwe\_hardness\_term}
  +\ec{module\_sis\_hardness\_term} \\
  &+\ec{bimodal\_to\_msis\_reduction\_loss\_term}
  +\ec{counted\_rom\_collision\_term} \\
  &+\ec{counted\_rom\_prequery\_term}
  +\ec{counted\_rom\_min\_entropy\_term}.
\end{aligned}
\]
Equivalently, under the concrete constants currently present in
\ec{HAETAE\_Reductions.ec}, the random-oracle programming contribution is
\(37/2^{58}\).

This implication does not include a checked equivalence
\(\mathsf{Spec}_{260204}\equiv\mathsf{ECModel}\).  Reading it as a theorem about
the February 2026 HAETAE specification therefore additionally relies on the
audited correspondences in \Cref{tab:internal-algorithm-fidelity-ledger} and on
the signing-randomness modeling bridge in
\Cref{assm:signing-randomness-modeling-bridge}.  The implication also does not
instantiate SHAKE in the standard model, derive concrete MLWE/Module-SIS
bit-security estimates, or check the running time of the reduction.  Those
items are boundary obligations recorded elsewhere in the dissertation.
\end{theorem}

\paragraph{Vacuity status.}
The checked implication above is accepted by EasyCrypt, but the current
final-hop antecedent
is not a non-vacuous concrete-security premise.  In
\ec{HAETAE\_Reductions.ec}, \ec{fs\_with\_aborts\_min\_entropy\_term} is \(12\)
and the checked facts imply
\(\ec{rejection\_sampling\_loss\_term}\ge 12\).  Since the right-hand side of
the displayed final-hop antecedent is a sum of two game-success probabilities,
it is at most \(2\).  The checked artifact therefore records the composition
that would follow if that hop premise were supplied; it should not be cited as a
useful concrete security theorem until the final-hop premise or the rejection
loss normalization is repaired.

\paragraph{Dependency form of the theorem.}
The checked implication has the logical shape
\[
  \mathsf{Checked}_{\easycrypt}
  \vdash
  \mathcal{H}_{\mathsf{EC}}
  \Longrightarrow
  \Adv_{\mathsf{ECModel},\rom}^{\eufcma}(A)
  \le \ec{haetae\_euf\_counted\_rom\_bound}.
\]
The specification-level reading has the strictly stronger, conditional shape
\[
  \mathsf{Checked}_{\easycrypt}
  \;+\;
  \mathcal{H}_{\mathsf{EC}}
  \;+\;
  \mathsf{Assumption}_{\mathsf{sign\text{-}coins}}
  \;+\;
  \mathsf{Audit}_{\mathsf{Spec}\leftrightarrow\mathsf{ECModel}}
  \Longrightarrow
  \Adv_{\mathsf{Spec}_{260204},\rom}^{\eufcma}(A)
  \lesssim
  \ec{haetae\_euf\_counted\_rom\_bound}.
\]
The symbol \(\lesssim\) is intentional: the dissertation does not supply a
machine-checked theorem equating the specification game with the EasyCrypt game.
The dependency form is the honest implication boundary and should be used whenever
the result is cited outside the formal-model setting.

The earlier broader ledger containing separate final summands for sampling,
rejection, extraction, and budget lifting is not the final counted-ROM theorem.
Those proof ideas are real, but in the checked counted theorem they are
internalized into the hop lemmas, the bimodal/MSIS reduction term, and the ROM
programming terms.  Treating them as additional final summands would double
count proof obligations that the EasyCrypt development has already composed.

\section{Hardness-term classification}
\label{sec:synthesis-hardness-classification}

The reduction file separates computational assumptions from proof-accounting
losses.  The classification is as follows.

\begin{longtable}{p{0.25\textwidth}p{0.30\textwidth}p{0.35\textwidth}}
\caption{Final counted-ROM bound terms.}
\label{tab:final-counted-rom-bound-terms}\\
\toprule
\textbf{Mathematical term} & \textbf{EasyCrypt identifier} & \textbf{Status in the proof} \\
\midrule
\endfirsthead
\toprule
\textbf{Mathematical term} & \textbf{EasyCrypt identifier} & \textbf{Status in the proof} \\
\midrule
\endhead
\(L_{\mathsf{MLWE}}\) & \ec{mlwe\_hardness\_term} & Abstract hardness term.  The
security theorem is conditional on this term being small for the relevant
parameter set. \\
\midrule
\(L_{\mathsf{ModSIS}}\) & \ec{module\_sis\_hardness\_term} & Abstract hardness term for the
Module-SIS layer. \\
\midrule
\(L_{\mathsf{bimodal}\to\mathsf{MSIS}}\) &
\ec{bimodal\_to\_msis\_reduction\_loss\_term} & Explicit reduction-loss slot,
currently defined as \(0\) in the EasyCrypt source. \\
\midrule
\(L_{\mathsf{coll}}\) & \ec{counted\_rom\_collision\_term} & Probability-accounting loss
for collisions among programmed or observed ROM sites. \\
\midrule
\(L_{\mathsf{pre}}\) & \ec{counted\_rom\_prequery\_term} & Counted-ROM
programmed-site prequery loss, with numerator
\(\ec{rom\_signature\_query\_budget}(\ec{rom\_hash\_query\_budget}+1)\). \\
\midrule
\(L_{\mathsf{me}}\) & \ec{counted\_rom\_min\_entropy\_term} & Counted-ROM
challenge min-entropy ledger term for the \ec{bad\_min\_entropy} predicate on
observed transcripts; distinct from sampler-prequery freshness. \\
\bottomrule
\end{longtable}

The table deliberately separates symbolic hardness terms from checked
probability-accounting terms.  The dissertation does not derive an independent
bit-security estimate for MLWE or Module-SIS at HAETAE levels 2, 3, and 5.
Those estimates belong to the cryptanalytic parameter analysis of the scheme.
The EasyCrypt theorem states that any such estimates may be inserted for
\ec{mlwe\_hardness\_term} and \ec{module\_sis\_hardness\_term}; it does not
prove the estimates themselves.

The counted-ROM prequery term should not be identified with the separate
\ec{sampler\_bad\_prequery} wrapper event.  The latter is set when a fresh
sampler-expansion query is already present either in the adversary hash log or
in the sampler self-log from earlier signing calls, and its checked budgeted
bound uses
\[
  \ec{rom\_signature\_query\_budget}
  \cdot
  \frac{\ec{rom\_hash\_query\_budget}+
        \ec{rom\_signature\_query\_budget}}
       {\ec{challenge\_support\_cardinality\_lower\_bound}}.
\]
That sampler-bad-prequery fact is part of the local active-signing and lifting
evidence; it is not the \ec{counted\_rom\_prequery\_term} summand in
\ec{haetae\_euf\_counted\_rom\_bound}.
Likewise, \ec{counted\_rom\_min\_entropy\_term} belongs to the counted-ROM
transcript ledger for challenge-entropy failures; it is not the support bound
for the sampler self-log event.

Consequently, the dissertation should not be read as a concrete NIST-level
security estimate.  It verifies the reduction ledger into symbolic MLWE and
Module-SIS assumptions.  A complete concrete-security chapter would additionally
need parameter-specific estimator evidence, assumption normalization, and a
clear rule for combining estimator costs with the checked ROM and reduction
losses.

The companion operation \ec{haetae\_euf\_bound} uses
\ec{rejection\_sampling\_loss\_term}; it is useful for paper-style abort and
reprogramming analysis.  The theorem emphasized here is instead the counted-ROM
variant \ec{haetae\_euf\_counted\_rom\_bound}.  The dissertation should not
conflate these two theorem shapes.

\section{Denotational reading of the proof scripts}
\label{sec:synthesis-denotational-final}

Each EasyCrypt procedure denotes a subprobability kernel.  If \(S\) is a state
space and \(R\) is the return space of a procedure \(c\), then
\[
  \llbracket c\rrbracket:S\to\mathcal{D}(S\times R).
\]
Losslessness lemmas assert that the kernel has total mass one.  An ordinary
Hoare judgment asserts partial correctness: every terminating outcome in the
support of the kernel satisfies the postcondition,
\[
  \{P\}\;c\;\{Q\}
  \quad\Longleftrightarrow\quad
  \forall s.\;P(s)\Rightarrow
  \operatorname{supp}(\llbracket c\rrbracket(s))
  \subseteq
  \{(s',r)\mid Q(s',r)\}.
\]
When \(c\) is additionally lossless, this support-inclusion statement can be
read as probability-one postcondition preservation:
\[
  \forall s.\;P(s)\Rightarrow
  \Pr_{(s',r)\leftarrow\llbracket c\rrbracket(s)}[Q(s',r)]=1.
\]
An equivalence judgment
\[
  \mathbf{equiv}[c_1\sim c_2:P\Rightarrow Q]
\]
asserts a relational coupling between the two kernels.  Thus an EasyCrypt game
hop is not merely a syntactic rewrite; it is a checked probabilistic relation
between two experiments.

The proof transitions have three mathematical forms:
\[
  \Pr[G_i:E_i]=\Pr[G_{i+1}:E_{i+1}],
\]
\[
  \Pr[G_i:E_i]\le \Pr[G_{i+1}:E_{i+1}]+\Pr[G_i:\mathsf{Bad}_i],
\]
and
\[
  \Pr[G_i:E_i]\le B_i(A).
\]
These correspond respectively to perfect couplings, fundamental-lemma
transitions, and direct reduction or probability bounds.

\section{Local clean-conditioned evidence and NMA-lifting gap}
\label{sec:synthesis-nonvacuity-final}

A formally checked lemma can still be cryptologically useless if its hypotheses
are never established.  The development contains local evidence against that
failure mode through concrete invariants over
\ec{ROMInternalTranscriptBudgetedPaperSimAsNMA.O.sign}.  The clean condition
contains the absence of a sampler prequery, sampler-ROM coverage, public-key
consistency, and current-secret-key consistency.  The top-level NMA path,
however, still uses some coarse fallback bounds proved from probability
boundedness and \(\ec{rejection\_sampling\_loss\_term}\ge 1\).  This section
therefore records local clean-conditioned evidence and the remaining
NMA-lifting gap, not a non-vacuous final-hop theorem.

The concrete lazy-ROM freshness laws are exposed by the following checked lemma
headers.

\begin{lstlisting}[language=EasyCrypt,caption={Concrete lazy-ROM sampler freshness laws.},label={lst:concrete-lazy-rom-freshness-laws}]
lemma concrete_lazy_rom_sampler_expand_query_fresh_le
    seed_coins (p : random_coins -> bool) &m :
  sampler_expand_query seed_coins \notin HAETAE_RO.FRO.m{m} =>
  Pr[HAETAE_RO.FRO.get(sampler_expand_query seed_coins) @ &m :
       p (ro_signing_coins res)] <=
  mu drandom_coins p.

lemma concrete_lazy_rom_sampler_expand_query_fresh_eq
    seed_coins (p : random_coins -> bool) :
  phoare[HAETAE_RO.FRO.get :
    arg = sampler_expand_query seed_coins /\
    sampler_expand_query seed_coins \notin HAETAE_RO.FRO.m
    ==> p (ro_signing_coins res)] =
  (mu drandom_coins p).
\end{lstlisting}

Mathematically, if \(r\leftarrow D_{coins}\) is fresh and
\(q=\mathsf{sampler\_expand\_query}(r)\) is not already in the lazy-ROM table,
then the distribution of \(\mathsf{ro\_signing\_coins}(H(q))\) is bounded by the
ideal coin distribution.  This is the exact statement needed to replace a bare
freshness assumption by a concrete lazy-ROM law.

The active signing invariant is recorded by the following theorem header.

\begin{lstlisting}[language=EasyCrypt,caption={Sampler-ROM coverage preservation for one active signing call.},label={lst:o-sign-clean-sampler-fresh}]
lemma concrete_budgeted_ro_signing_attempt_o_sign_clean_sampler_fresh :
  hoare[ROMInternalTranscriptBudgetedPaperSimAsNMA
          (A, ROSigningAttemptPaperSimSampler(HAETAE_RO.FRO),
           HAETAE_RO.FRO).O.sign :
    sampler_rom_covered
      HAETAE_RO.FRO.m
      ROMInternalTranscriptBudgetedPaperSimAsNMA.adversary_hash_queries
      ROMInternalTranscriptBudgetedPaperSimAsNMA.sampler_expand_queries /\
    ! ROMInternalTranscriptBudgetedPaperSimAsNMA.sampler_bad_prequery ==>
    sampler_rom_covered
      HAETAE_RO.FRO.m
      ROMInternalTranscriptBudgetedPaperSimAsNMA.adversary_hash_queries
      ROMInternalTranscriptBudgetedPaperSimAsNMA.sampler_expand_queries].
\end{lstlisting}

The one-call loss theorem is then applied through the concrete surface lemma
\ec{concrete\_budgeted\_o\_sign\_clean\_one\_call\_loss\_from\_self\_logged\_surface},
whose conclusion has the form
\[
\begin{aligned}
&\Pr[O_{\mathsf{attempt}}.\mathsf{sign}(m,ctx):
      p(res)\wedge\neg\mathsf{sampler\_bad\_prequery}] \\
&\qquad\le
  \Pr[O_{\mathsf{exact}}.\mathsf{sign}(m,ctx):p(res)]
  + \ec{rejection\_sampling\_loss\_term}.
\end{aligned}
\]
This displayed conclusion is conditional on the checked premises in the EasyCrypt
lemma: the structural sample-loss obligation, sampler-ROM coverage, absence of
the current sampler-prequery event, an active signing budget, current public-key
and secret-key consistency, and the attempt-to-sample projection premise.
This is the local theorem surface that a non-vacuous active-signing lift should
consume.  The present top-level NMA lemmas do not yet fully use this surface as
their final proof path; some later bridge lemmas still fall back to coarse
probability bounds whose usefulness depends on the oversized
\ec{rejection\_sampling\_loss\_term}.  Future maintenance should replace those
fallbacks by a final NMA-level composition that consumes the clean one-call
surface directly.

\section{Wrapper counts and runtime cost}
\label{sec:synthesis-budget-runtime-final}

The checked implication uses a counted-ROM bound ledger and budgeted wrapper
lemmas with fixed hash-query and signing-query wrapper constants.  In the
current reduction file these constants are fixed to
\[
  q_H=2,\qquad q_S=2.
\]
These constants account for the wrapper layer; they are not syntactic query caps
on the source \ec{Sig\_ROM.EUF\_CMA} game.  The budgeted wrapper proof exposes
the state needed to lift one-call clean
losses through adaptive oracle interaction without replacing the proof by an
unstructured union-bound envelope.  The remaining obligation is to make the
final NMA-level composition use that structure directly instead of relying on
coarse fallback bounds.

Running time is different.  The EasyCrypt theorem checked here does not contain
a cost semantics for the reduction.  The paper-level runtime statement should
therefore be written as an external interpretation ledger:
\[
  \operatorname{Time}(B_A)
  \le
  \operatorname{Time}(A)
  +T_{\mathsf{ROM}}(q_H)
  +T_{\mathsf{sign}}(q_S)
  +T_{\mathsf{extract}}(q_H,q_S).
\]
Here \(T_{\mathsf{ROM}}\) is lazy-table maintenance, \(T_{\mathsf{sign}}\) is
simulated signing and sampler instrumentation, and \(T_{\mathsf{extract}}\) is
the final reduction/extraction overhead.  This ledger is not a checked theorem
unless a future cost-aware EasyCrypt development adds explicit cost judgments.
Consequently, the formal theorem statement in this dissertation should be read
as a probability theorem with query-budget accounting, not as a theorem about
the concrete running time of the declared reduction modules.

\section{Trusted base}
\label{sec:synthesis-trusted-base-final}

The trusted base has five parts:

\begin{enumerate}[label=\arabic*.]
\item EasyCrypt's logic, libraries, module system, and probability semantics.
\item The manifest discipline selecting the files in \ec{provable-security/proof-files.txt}.
\item The interpretation of \ec{mlwe\_hardness\_term} and
\ec{module\_sis\_hardness\_term} as external lattice-hardness assumptions.
\item The model/spec correspondence between \ec{HAETAE\_v260204.pdf} and the
formal scheme model.
\item The fact that implementation-level byte parsing, arithmetic, hashing, and
constant-time behavior are outside the central provable-security theorem.
\end{enumerate}

This is the precise scope of the dissertation claim: the formal HAETAE ROM
proof-surface implication is machine-checked, and the remaining assumptions are
stated at the boundary rather than hidden inside the proof narrative.
